\documentclass[a4paper,10pt]{article}
\title{A Metrically Exact Article}
\author{Kieran Meinhardt}
\date{September 5, 2026}
\begin{document}
\maketitle

\begin{abstract}
Page geometry, baseline grid, fonts and sectioning are taken from the class
files themselves, and checked against real output.
\end{abstract}

\section{Running text}

Body text is set in Computer Modern at ten points on a twelve point baseline
grid.\footnote{The footnote rule and type size come from the class files too.}

\subsection{Lists}

The itemize and enumerate environments use the class's own list parameters:

\begin{itemize}
\item The item body sits two and a half ems in.
\item Labels are set flush right, ending half an em earlier.
\end{itemize}

\begin{enumerate}
\item Items are separated by itemsep plus parsep.
\item The list is surrounded by topsep plus partopsep.
\end{enumerate}

\subsubsection{Displayed material}

A displayed equation gets the class's display skips:
\[ e^{i\pi} + 1 = 0 \]
and the text after it stays on the grid.

\paragraph{Run-in headings} These are set with a negative afterskip, so the
heading runs into the paragraph rather than standing on its own line.

\section{Floats}

\begin{figure}[h]
\centering
\rule{4cm}{1.4cm}
\caption{A caption short enough to be centred.}
\end{figure}

\end{document}
